\subsection{THE FIELD OF RATIONAL NUMBERS}

DEFINITION 4. The field of fractions of the ring Z of rational integers is called the field of rational numbers and is denoted by Q. The elements of Q are called rational numbers.

Every rational number is thus of the form \(a / b\) where \(a\) and \(b\) are rational integers with \(b \neq 0\) (and we may even take \(b > 0\) as the relation

\[
a / b = (- a) / (- b)
\]

proves). \(\mathbf{Q}_{+}\) is used to denote the set of rational numbers of the form \(a / b\) with \(a\in \mathbf{N}\) and \(b\in \mathbf{N}^*\) .

We have the relations:

\begin{enumerate}
\def\labelenumi{(\arabic{enumi})}
\tightlist
\item
\end{enumerate}

\[
\mathbf {Q} _ {+} + \mathbf {Q} _ {+} = \mathbf {Q} _ {+}\tag{2}
\]

\[
\mathbf {Q} _ {+}. \mathbf {Q} _ {+} = \mathbf {Q} _ {+}\tag{3}
\]

\[
\mathbf {Q} _ {+} \cap (- \mathbf {Q} _ {+}) = \{0 \}\tag{4}
\]

\[
\mathbf {Q} _ {+} \cup (- \mathbf {Q} _ {+}) = \mathbf {Q}\tag{5}
\]

\[
\mathbf {Q} _ {+} \cap \mathbf {Z} = \mathbf {N}.
\]

The first two follow from the formulae \(a / b + a' / b' = (ab' + a'b) / bb'\) , \((a / b)(a'/b') = aa'/bb'\) , \(0 \in \mathbf{Q}_+\) , \(1 \in \mathbf{Q}_+\) and the fact that \(\mathbf{N}\) is stable under addition and multiplication and \(\mathbf{N}^*\) under multiplication. Obviously \(0 \in \mathbf{Q}_+\) , whence \(0 \in (-\mathbf{Q}_+)\) ; let \(x\) be in \(\mathbf{Q}_+ \cap (-\mathbf{Q}_+)\) ; then there exist positive integers \(a, b, a', b'\) with \(b \neq 0\) , \(b' \neq 0\) and \(x = a / b = -a'/b'\) ; then \(ab' + ba' = 0\) , whence \(ab' = 0\) (for \(ab' \geqslant 0\) and \(ba' \geqslant 0\) ) and therefore \(a = 0\) since \(b' \neq 0\) ; in other words, \(x = 0\) . Finally, obviously \(\mathbf{N} \subset \mathbf{Z}\) and \(\mathbf{N} \subset \mathbf{Q}_+\) . Conversely, if \(x\) belongs to \(\mathbf{Z} \cap \mathbf{Q}_+\) , it is a rational integer; there exist two rational integers \(a\) and \(b\) with \(a \geqslant 0\) , \(b > 0\) and \(x = a / b\) , whence \(a = bx\) ; if \(x \notin \mathbf{N}\) , then \(-x > 0\) , whence \(-a = b(-x) > 0\) and consequently \(a < 0\) contrary to the hypothesis.

Given two rational numbers \(x\) and \(y\) , we write \(x \leqslant y\) if \(y - x \in \mathbf{Q}_+\) . It is easily deduced from formulae (1), (3) and (4) that \(x \leqslant y\) is a total ordering on

Q, from (5) that this relation induces the usual order relation on Z. Finally, from (1) it follows that the relations \(x \leqslant y\) and \(x' \leqslant y'\) imply \(x + x' \leqslant y + y'\) and from (2) that the relation \(x \leqslant y\) implies \(xz \leqslant yz\) for all \(z \geqslant 0\) and \(xz \geqslant yz\) for \(z \leqslant 0\) (cf.~VI, § 2, no. 1).

Let \(x\) be a rational number. \(x\) is said to be positive if \(x \geqslant 0\) , strictly positive if \(x > 0\) , negative if \(x \leqslant 0\) and strictly negative if \(x < 0\) . The set of positive rational numbers is \(\mathbf{Q}_{+}\) and that of negative numbers is \(-\mathbf{Q}_{+}\) . If \(\mathbf{Q}^{*}\) denotes the set of non-zero rational numbers, the set \(\mathbf{Q}_{+}^{*}\) of strictly positive numbers is equal to \(\mathbf{Q}^{*} \cap \mathbf{Q}_{+}\) and \(-\mathbf{Q}_{+}^{*}\) is the set of strictly negative numbers.

The sets \(\mathbf{Q}_{+}^{*}\) and \(\{1, -1\}\) are subgroups of the multiplicative group \(\mathbf{Q}^{*}\) . Every rational number \(x \neq 0\) can be expressed in one and only one way in the form \(1.y\) , \((-1).y\) , where \(y > 0\) ; hence the multiplicative group \(\mathbf{Q}^{*}\) is the product of the subgroups \(\mathbf{Q}_{+}^{*}\) and \(\{1, -1\}\) , the component of \(x\) in \(\mathbf{Q}_{+}^{*}\) is called the absolute value of \(x\) and is denoted by \(|x|\) ; the component of \(x\) in \(\{-1, 1\}\) (equal to 1 if \(x > 0\) , to \(-1\) if \(x < 0\) ) is called the sign of \(x\) and denoted by \(\operatorname{sgn} x\) .

Usually these two functions are extended to the whole of \(\mathbf{Q}\) by setting \(|0| = 0\) and \(\operatorname{sgn} 0 = 0\) .

